\section{Experiments}
\label{sec:experiments}

We evaluate two claims: (1)~H1---that forgetting is predicted by drift in importance-ranked activation dimensions (\cref{sec:h1_validation}), and (2)~that SFP's causal mechanism is sound---perturbing important dimensions causes disproportionate performance degradation (\cref{sec:causal}).
We also report an $r$-sweep diagnostic that characterizes the effective dimensionality of forgetting (\cref{sec:rsweep}).

\subsection{Experimental Setup}
\label{sec:setup}

\paragraph{Models.}
We evaluate on two scales: SmolLM2-135M-Instruct (135M parameters, CPU) and Qwen2.5-1.5B-Instruct (1.5B parameters, single GPU).
Both models are fine-tuned with LoRA~\citep{hu2022lora} (rank 32, $\alpha = 64$, applied to all linear layers).

\paragraph{Tasks.}
Our primary continual learning setting is Math$\to$Code: the model is first evaluated on GSM8K~\citep{cobbe2021training} (math reasoning), then fine-tuned on MBPP~\citep{austin2021program} (code generation).
Forgetting is measured as the accuracy drop on GSM8K after fine-tuning on MBPP.
The full benchmark (\cref{sec:full_benchmark}) extends to five task pairs covering instruction following (Alpaca$\to$IFEval), safety (HH-RLHF), and domain adaptation (PubMed).

\paragraph{SFP configuration.}
We extract activations at three probe layers ($L/4$, $L/2$, $3L/4$), compute PCA with $k = 128$ components, and preserve the top $r = 32$ dimensions ranked by ablation importance.
The memory set consists of 128 examples from the old task.
Training runs for 500 steps with checkpoints every 25 steps, yielding 20 evaluation points for regression analysis.

\paragraph{Metrics.}
For H1 validation, we compute $R^2$ of the linear regression of forgetting (accuracy drop) on projected drift $\| P_W(\Delta a) \|$ across checkpoints.
We compare four predictors: top-$r$ (ablation-ranked), total (all dimensions), random (uniformly sampled $r$ dimensions), and bottom-$r$ (least important by ablation).
For the causal test, we measure the perplexity increase ($\Delta$ppl) when injecting noise into each set of dimensions at inference time.

\paragraph{Baselines.}
We compare against: (1)~\textbf{Naive} fine-tuning with no forgetting mitigation; (2)~\textbf{Replay} with interleaved old-task examples (same 128-example budget as SFP's memory set); (3)~\textbf{LwF}~\citep{li2017learning}: logit distillation from the pre-trained model; (4)~\textbf{Hidden distillation}: full hidden-state distillation in the style of PODNet~\citep{douillard2020podnet}; (5)~\textbf{Orthogonal LoRA}~\citep{qin2024exploring}: constraining adapter updates to be orthogonal to previous task directions.
For SFP, we additionally report two ablation controls: \textbf{SFP-random} (preserving $r$ random PCA dimensions instead of ablation-ranked ones) and \textbf{SFP-random-basis} (preserving $r$ random directions not from PCA).

\subsection{H1 Validation: Correlational Evidence}
\label{sec:h1_validation}

\paragraph{1.5B results (Qwen2.5-1.5B-Instruct).}
\Cref{tab:h1_results} reports the $R^2$ values for predicting forgetting from activation drift under different predictor subspaces.
The top-$r$ ablation-ranked dimensions achieve $R^2 = 0.933$, outperforming total drift ($0.888$), random dimensions ($0.850$), and bottom-ranked dimensions ($0.827$).
The ranking is exactly as predicted by H1: top $>$ total $>$ random $>$ bottom.

\begin{table}[t]
\centering
\caption{$R^2$ of linear regression predicting forgetting from activation drift at 1.5B scale (Qwen2.5-1.5B-Instruct, Math$\to$Code, 500 steps, 20 checkpoints). Higher is better. H1 predicts: top $>$ total $>$ random $>$ bottom.}
\label{tab:h1_results}
\begin{tabular}{lcc}
\toprule
Predictor & Dims & $R^2$ \\
\midrule
Top-$r$ (ablation) & 32 & \textbf{0.933} \\
Total drift & 128 & 0.888 \\
Random dims & 32 & 0.850 \\
Bottom-$r$ (ablation) & 32 & 0.827 \\
\bottomrule
\end{tabular}
\end{table}

Against our pre-registered gate criteria, 2 of 3 checks pass: (i)~$R^2(\text{top-}r) = 0.933 > 0.6$~\checkmark; (ii)~$R^2(\text{top-}r) > R^2(\text{total})$~\checkmark; (iii)~$R^2(\text{top-}r) \gg R^2(\text{random})$---the margin is $0.083$, which we judge as modest rather than decisive~($\times$).
We score this as a \textbf{partial pass}: the correct ranking emerges, but the margin over random is not overwhelming.

\paragraph{135M results (SmolLM2-135M-Instruct).}
At the 135M scale, H1 \textbf{fails}: the $R^2$ ranking is inverted, with bottom-ranked dimensions ($R^2 = 0.89$) outperforming top-ranked ($R^2 = 0.71$).
We attribute this to a scale artifact: at 135M parameters, the model's activation space may be too compressed for meaningful separation between important and unimportant feature dimensions.
This negative result motivates our focus on 1.5B+ models and highlights that H1 is a scale-dependent hypothesis.

\paragraph{Baseline comparison.}
The baseline math accuracy before fine-tuning is $44.44\%$ on GSM8K for Qwen2.5-1.5B-Instruct.
\todo{Full baseline comparison table (SFP vs.\ Naive, Replay, LwF, Hidden, OrthoLoRA) pending GPU benchmark runs.}

\subsection{H1 Validation: Causal Evidence}
\label{sec:causal}

The correlational analysis in \cref{sec:h1_validation} establishes that top-$r$ drift \emph{predicts} forgetting, but correlation does not imply causation.
To test the causal direction, we inject calibrated noise into specific activation dimensions at inference time and measure the effect on old-task perplexity.
If the top-$r$ dimensions are causally important for old-task performance, perturbing them should cause more degradation than perturbing random or bottom-ranked dimensions.

\paragraph{Protocol.}
For each set of dimensions (top-$r$, random, bottom-$r$), we add Gaussian noise $\epsilon \sim \mathcal{N}(0, \sigma^2 I)$ to the selected PCA components of the activation at each probe layer during a forward pass on old-task evaluation data.
The noise scale $\sigma$ is calibrated to the standard deviation of the projected activations, so the perturbation is comparable in magnitude across dimension sets.

\paragraph{Results (135M, CPU smoke test).}
Even at the 135M scale where correlational H1 fails, the causal test passes clearly (\cref{tab:causal}):

\begin{table}[t]
\centering
\caption{Causal intervention test: perplexity increase ($\Delta$ppl) when injecting noise into different activation dimension sets (SmolLM2-135M, Math task). Higher $\Delta$ppl indicates greater causal importance.}
\label{tab:causal}
\begin{tabular}{lc}
\toprule
Dimension Set & $\Delta$ppl \\
\midrule
Top-$r$ (ablation) & $+0.85$ \\
Random & $+0.21$ \\
Bottom-$r$ (ablation) & $+0.02$ \\
\bottomrule
\end{tabular}
\end{table}

Perturbing the top-$r$ dimensions causes $4\times$ more degradation than random dimensions and $42\times$ more than bottom-ranked dimensions.
This result is striking because it holds at 135M parameters, where the correlational H1 test fails.
The discrepancy suggests that ablation importance correctly identifies causally important dimensions even at small scale, but the \emph{correlation} between drift in those dimensions and forgetting is confounded at 135M---possibly because the model's limited capacity means drift in ``unimportant'' dimensions also contributes to forgetting.

\paragraph{1.5B causal test.}
\todo{Results from the 1.5B GPU causal intervention run are pending. We expect the causal signal to be at least as strong as at 135M, given the stronger correlational evidence at that scale.}

\subsection{Dimensionality Sweep}
\label{sec:rsweep}

To characterize the effective dimensionality of forgetting, we sweep $r \in \{1, 4, 32, 128\}$ and measure both the variance explained by the top-$r$ PCA dimensions and the $R^2$ of forgetting prediction (\cref{tab:rsweep}).

\begin{table}[t]
\centering
\caption{$r$-sweep at 1.5B scale: variance explained by top-$r$ PCA dimensions and forgetting prediction $R^2$. Even $r = 1$ captures 96.2\% of variance and achieves $R^2 = 0.854$.}
\label{tab:rsweep}
\begin{tabular}{rcc}
\toprule
$r$ & Variance Explained (\%) & Forgetting $R^2$ \\
\midrule
1   & 96.2 & 0.854 \\
4   & 96.9 & 0.871 \\
32  & 98.7 & 0.882 \\
128 & 100.0 & 0.884 \\
\bottomrule
\end{tabular}
\end{table}

Two observations stand out.
First, a single PCA dimension captures 96.2\% of the activation variance, consistent with the highly concentrated spectrum typical of transformer activations.
Second, forgetting $R^2$ increases only modestly from $r = 1$ ($0.854$) to $r = 128$ ($0.884$), suggesting that the forgetting-relevant signal is concentrated in very few dimensions.
However, we note that the ablation-ranked $R^2$ in \cref{tab:h1_results} ($0.933$ for $r = 32$) is higher than the PCA-ranked $R^2$ here ($0.882$), confirming that ablation importance provides a better ranking than variance explained.

\subsection{Full Benchmark}
\label{sec:full_benchmark}

\todo{Full five-task benchmark results (Math, Code, Instruction Following, Safety, Domain) with all baselines are pending GPU runs. The benchmark measures Average Accuracy (AA), Backward Transfer (BWT), Forward Transfer (FWT), and per-task forgetting across both model scales.}
